\documentclass[conference]{IEEEtran}
\IEEEoverridecommandlockouts
\usepackage{times}
\usepackage{cite}
\usepackage{amsmath,amssymb,amsfonts}
\usepackage{graphicx}
\usepackage{textcomp}
\usepackage{xcolor}
\usepackage[hidelinks]{hyperref}
\usepackage{enumitem}
\setlist{nosep, leftmargin=*}

\def\BibTeX{{\rm B\kern-.05em{\sc i\kern-.025em b}\kern-.08em
    T\kern-.1667em\lower.7ex\hbox{E}\kern-.125emX}}

\begin{document}

\title{El Nivel de Transporte en Redes de Computadoras: Análisis Comparativo de TCP/UDP y Factores de Desempeño}

\author{
\IEEEauthorblockN{Jorge Moreno}
\IEEEauthorblockA{Licenciatura en Computación, Séptimo Semestre\\
Redes de Computadoras 2\\
Correo:  jmoreno29@uc.edu.ve}
}

\maketitle

\begin{abstract}
La capa de transporte constituye el eslabón que permite la comunicación lógica extremo a extremo entre procesos de aplicación, independientemente de la infraestructura de red subyacente. Este artículo revisa de manera integral los fundamentos del nivel de transporte, contrastando el modelo no orientado a conexión (UDP) con el modelo orientado a conexión (TCP), incluyendo mecanismos de multiplexación, transferencia fiable de datos, control de flujo y control de congestión. Asimismo, se examinan protocolos auxiliares como RPC y RTP, y se analizan los principales factores que determinan el desempeño de una red, junto con estrategias de diseño orientadas a optimizar el procesamiento de las unidades de datos del protocolo de transporte (TPDU). Se concluye con criterios prácticos para seleccionar el protocolo de transporte adecuado según los requisitos de la aplicación.
\end{abstract}

\begin{IEEEkeywords}
capa de transporte, TCP, UDP, sockets, control de congestión, control de flujo, desempeño de redes, TPDU, RPC, RTP
\end{IEEEkeywords}

\section{Introducción y servicios de transporte}

La capa de transporte es la cuarta capa del modelo de referencia OSI y la primera orientada a la aplicación cuya comunicación con su entidad par es estrictamente extremo a extremo \cite{tanenbaum}. A diferencia de las capas inferiores, que intercambian unidades de datos punto a punto con el siguiente elemento de la red, el protocolo de transporte se ejecuta únicamente en los sistemas finales (\textit{hosts}), nunca en los encaminadores intermedios \cite{kurose}. En la pila TCP/IP, que es la que rige de facto en Internet, esta capa concentra las funciones que en el modelo OSI de siete niveles se reparten entre las capas de transporte y sesión, y define dos protocolos principales —UDP y TCP— con modelos de servicio radicalmente distintos entre sí, además de alternativas más recientes como SCTP y DCCP, orientadas a escenarios híbridos que requieren, por ejemplo, entrega de mensajes fuera de orden con control de flujo, o control de congestión sin garantía de entrega \cite{fuoc}. La elección del protocolo adecuado no es un detalle de implementación menor: condiciona directamente la latencia percibida, la resiliencia ante pérdidas y la escalabilidad del sistema distribuido que se construye sobre ella.

\subsection{Servicios ofrecidos}

La capa de transporte proporciona una \emph{comunicación lógica} entre procesos de aplicación ejecutándose en máquinas distintas: desde la perspectiva de la aplicación, ambos extremos parecen conectados directamente, aunque en realidad la comunicación atraviesa múltiples encaminadores y enlaces heterogéneos \cite{fuoc}. Entre los servicios característicos se destacan:

\begin{itemize}
    \item \textbf{Segmentación y reensamblaje}: los mensajes de aplicación se fragmentan en segmentos (TPDU o PDU-4) antes de pasarlos a la capa de red, y se reconstruyen en el receptor.
    \item \textbf{Multiplexación/demultiplexación}: permite que múltiples procesos comparen una única conexión de red, mediante el uso de \textit{sockets} identificados por puertos.
    \item \textbf{Control de flujo}: regula la velocidad de envío para no saturar al receptor.
    \item \textbf{Control de errores y entrega fiable} (opcional, según el protocolo): garantiza integridad, orden y ausencia de duplicados mediante confirmaciones (ACK).
    \item \textbf{Control de congestión}: dirigido a la red, no a la aplicación, para evitar el colapso de los recursos compartidos.
\end{itemize}

Cabe destacar que no todos los protocolos de transporte ofrecen la totalidad de estos servicios; UDP, por ejemplo, renuncia deliberadamente a la fiabilidad y al control de flujo en favor de la simplicidad y la baja latencia \cite{fuoc}.

\subsection{Control de la conexión y primitivas de transporte}

Los protocolos orientados a conexión atraviesan tres fases: \emph{establecimiento}, \emph{transferencia de datos} y \emph{liberación}. Las \textit{primitivas del servicio de transporte} (por ejemplo, \texttt{CONNECT}, \texttt{SEND}, \texttt{RECEIVE}, \texttt{DISCONNECT}) constituyen la interfaz que exponen los protocolos hacia las capas superiores, permitiendo a la aplicación invocar estas fases sin conocer los detalles del protocolo subyacente.

Para el direccionamiento, cada extremo necesita un identificador único, llamado \textit{TSAP} en OSI, que en Internet se conoce como el número de puerto. Un \textit{socket} queda definido por (IP, puerto), y en TCP por la 4-tupla (IP origen, puerto origen, IP destino, puerto destino), lo que permite demultiplexar segmentos hacia distintos procesos incluso cuando comparten el mismo puerto destino \cite{kurose}.

\subsection{Sockets de Berkeley}

La interfaz de \textit{sockets} de Berkeley constituye la implementación práctica de estas primitivas en sistemas Unix y compatibles. Un socket es un punto terminal de comunicación identificado por una dirección IP y un puerto; las operaciones principales son \texttt{socket()}, \texttt{bind()}, \texttt{listen()}, \texttt{accept()}, \texttt{connect()}, \texttt{send()}/\texttt{write()} y \texttt{recv()}/\texttt{read()}. Existen sockets de tipo \texttt{SOCK\_STREAM} (orientados a conexión, sobre TCP) y \texttt{SOCK\_DGRAM} (no orientados a conexión, sobre UDP), cada uno adecuado a distintas necesidades de fiabilidad y latencia.\\

\begin{table}[htbp]
\caption{Principales campos de la cabecera TCP (RFC 793)}
\label{tab:tcpheader}
\centering
\begin{tabular}{|p{2.3cm}|p{4.7cm}|}
\hline
\textbf{Campo} & \textbf{Función} \\
\hline
Puerto origen/destino & Identifican los procesos de aplicación (demultiplexación) \\
\hline
Núm. de secuencia & Posición, en bytes, del primer byte del segmento \\
\hline
Núm. de reconocimiento & Próximo byte esperado por el receptor (ACK) \\
\hline
Indicadores (SYN, ACK, FIN, RST, PSH, URG) & Controlan el ciclo de vida y semántica del segmento \\
\hline
Ventana anunciada & Espacio disponible en el \textit{buffer} de recepción \\
\hline
Suma de comprobación & Detección de errores (obligatoria, a diferencia de UDP) \\
\hline
Opciones (MSS, \textit{window scale}, SACK, \textit{timestamp}) & Negociación de parámetros de rendimiento \\
\hline
\end{tabular}
\end{table}

\subsection{Elementos clave de los protocolos de transporte}

Independientemente del protocolo concreto, los elementos comunes de diseño incluyen: direccionamiento, establecimiento y liberación de la conexión, control de flujo mediante memorias intermedias (\textit{buffers}) dinámicas, y multiplexación/demultiplexación. El establecimiento de una conexión es particularmente delicado en redes donde los paquetes pueden perderse, duplicarse o retrasarse; la solución adoptada por TCP es el acuerdo en tres vías (\textit{three-way handshake}), que evita que segmentos duplicados de conexiones antiguas sean interpretados erróneamente como pertenecientes a una nueva conexión \cite{tanenbaum}.

\section{Protocolos sin conexión: UDP}

\subsection{Características principales}

El \emph{User Datagram Protocol} (UDP, RFC~768) es un protocolo de transporte no orientado a conexión que constituye, en esencia, una extensión mínima de IP con capacidad de multiplexación por puertos \cite{fuoc}. Sus características esenciales son:

\begin{itemize}
    \item No existe fase de establecimiento de conexión, por lo que no introduce retardos adicionales (ejemplo típico: DNS).
    \item No mantiene estado de la conexión en el sistema final.
    \item Cabecera mínima de 8 bytes (puerto origen, puerto destino, longitud y \textit{checksum}), frente a los 20 bytes mínimos de TCP.
    \item No hay confirmaciones (ACK), ni control de flujo, ni control de congestión: cada datagrama se trata de forma independiente.
    \item Ofrece únicamente detección de errores mediante suma de comprobación (opcional); si se detecta un error, el datagrama se descarta sin notificar al emisor.
\end{itemize}

La ausencia de estos mecanismos traslada toda la responsabilidad de fiabilidad, orden y control de duplicados a la aplicación. Esto resulta ventajoso en aplicaciones de tiempo real —telefonía IP, videoconferencia, monitoreo de señales— donde un dato retrasado es peor que un dato perdido, ya que introducir retransmisiones generaría \textit{jitter} inaceptable \cite{fuoc}.

\begin{table}[htbp]
\caption{Comparación entre UDP y TCP}
\label{tab:comparacion}
\centering
\begin{tabular}{|p{1.6cm}|p{2.7cm}|p{2.7cm}|}
\hline
\textbf{Aspecto} & \textbf{UDP} & \textbf{TCP} \\
\hline
Conexión & No orientado & Orientado (3-way handshake) \\
\hline
Fiabilidad & No garantizada & Garantizada (ACK, retransmisión) \\
\hline
Orden & No garantizado & Garantizado (num. secuencia) \\
\hline
Control de flujo & Ausente & Ventana deslizante \\
\hline
Control de congestión & Ausente & \textit{Slow start}, AIMD \\
\hline
Cabecera & 8 bytes & 20--60 bytes \\
\hline
Unidad & Datagrama & Flujo de bytes \\
\hline
Uso típico & DNS, RTP, RPC, VoIP & HTTP, FTP, SMTP, SSH \\
\hline
\end{tabular}
\end{table}


\subsection{Llamada a Procedimiento Remoto (RPC)}

El modelo de \emph{Remote Procedure Call} (RPC) permite invocar procedimientos ejecutados en una máquina remota como si fueran llamadas locales, ocultando los detalles del transporte subyacente. Un cliente invoca una función; un \textit{stub} local serializa (\textit{marshalling}) los parámetros y los envía al servidor, típicamente sobre UDP dado el patrón de comunicación de solicitud--respuesta de corta duración. El servidor ejecuta el procedimiento y devuelve el resultado por el mismo mecanismo. Cuando UDP se emplea como transporte de RPC, es responsabilidad de la propia capa de RPC implementar mecanismos ligeros de reintento y detección de duplicados (mediante identificadores de transacción), ya que UDP no los provee.

\subsection{Protocolo de Transporte en Tiempo Real (RTP)}

El \emph{Real-time Transport Protocol} (RTP, RFC~3550) se ejecuta habitualmente sobre UDP y añade a los datagramas información necesaria para la reproducción de flujos de audio y video: número de secuencia (para detectar pérdidas y reordenar), marca de tiempo (\textit{timestamp}, para la sincronización de reproducción) y un identificador de la fuente. RTP no proporciona por sí mismo fiabilidad ni control de congestión; estas funciones se delegan al protocolo compañero \emph{RTCP} (\textit{RTP Control Protocol}), encargado de reportar estadísticas de calidad de servicio (pérdida de paquetes, \textit{jitter}) entre los participantes de la sesión. La combinación UDP+RTP resulta idónea para aplicaciones multimedia donde la puntualidad de entrega prevalece sobre la fiabilidad absoluta.

\end{document}
